\section*{Capítulo \ref{ch:g10:exercise-and-energy} --- El ejercicio y las necesidades energéticas del cuerpo}

\begin{solution}{exo:g10:exercise-and-energy:1}
La energía gastada en reposo completo para mantener el cuerpo vivo ---
corazón, cerebro, temperatura, mantenimiento celular: unos \qty{80}{W},
unos \qty{7000}{kJ} al día.
\end{solution}

\begin{solution}{exo:g10:exercise-and-energy:2}
$2.0 \times 20 = \qty{40}{kJ}$ por minuto, es decir $40000/60 \approx
\qty{670}{W}$.
\end{solution}

\begin{solution}{exo:g10:exercise-and-energy:3}
La reserva fosforada del músculo (unos diez segundos); la \omterm{def:g10:cell-metabolism:fermentation}{fermentación}
láctica (uno o dos minutos de esfuerzo máximo); la \omterm{prop:g10:cell-metabolism:respiration}{respiración celular}
(horas).
\end{solution}

\begin{solution}{exo:g10:exercise-and-energy:4}
La mayor velocidad a la que el cuerpo puede captar y usar oxígeno.
$3500/70 = \qty{50}{mL/min}$ por kilogramo.
\end{solution}

\begin{solution}{exo:g10:exercise-and-energy:5}
En los músculos, unos \qty{400}{g}, que alimentan al músculo que lo
almacena; en el hígado, unos \qty{100}{g}, liberados como glucosa a la
sangre para el cerebro y los demás tejidos.
\end{solution}

\begin{solution}{exo:g10:exercise-and-energy:6}
\qty{2.7}{L/min}: $2.7 \times 20 = \qty{54}{kJ}$ por minuto, es decir
\qty{900}{W}. Calor: $900 - 200 = \qty{700}{W}$.
\end{solution}

\begin{solution}{exo:g10:exercise-and-energy:7}
Trabajo $mgh = 60 \times 9.8 \times 800 \approx \qty{470}{kJ}$. Con un
rendimiento del 25\%, los músculos gastaron unos \qty{1900}{kJ}. En
\qty{7200}{s}: unos \qty{260}{W} por encima del reposo.
\end{solution}

\begin{solution}{exo:g10:exercise-and-energy:8}
Un esprint de 100 metros va casi por completo con la reserva fosforada:
sin deuda de oxígeno y con poco ácido láctico. Una carrera de 400
metros dura alrededor de un minuto y saca la mitad de su energía de la
\omterm{def:g10:cell-metabolism:fermentation}{fermentación} láctica: el ácido láctico se acumula (de ahí el dolor) y
la respiración que debe eliminarlo mantiene al corredor jadeando
durante minutos.
\end{solution}

\begin{solution}{exo:g10:exercise-and-energy:9}
Energía de los \omterm{def:g10:chemistry-of-life:families}{glúcidos} $0.6 \times 3000 = \qty{1800}{kJ}$; \omterm{def:g10:exercise-and-energy:glycogen}{glucógeno}
$1800/17 \approx \qty{106}{g}$; agua liberada, unos \qty{320}{g}.
\end{solution}

\begin{solution}{exo:g10:exercise-and-energy:10}
Potencia metabólica de unos \qty{600}{W}: $\num{1.1e6}/600 \approx
\qty{1830}{s}$, unos 30 minutos.
\end{solution}

\begin{solution}{exo:g10:exercise-and-energy:11}
Tres cuartas partes de la energía liberada en los músculos son calor,
varios cientos de vatios durante un ejercicio duro. El sudor lo
elimina: evaporar agua cuesta unos \qty{2.4}{kJ} por gramo, y el flujo
de sangre hacia la piel lleva el calor hasta allí. Sin sudor, la
temperatura corporal subiría un grado cada pocos minutos.
\end{solution}

\begin{solution}{exo:g10:exercise-and-energy:12}
Cuesta arriba, la potencia necesaria es proporcional a la masa
corporal: el sujeto de \qty{60}{kg} tiene \qty{60}{mL/min} por
kilogramo frente a \num{40}, y sube más deprisa. En la bicicleta de
laboratorio la potencia no depende de la masa, ambos tienen los mismos
\qty{3.6}{L/min} y ambos pueden sostener los mismos vatios.
\end{solution}

\begin{solution}{exo:g10:exercise-and-energy:13}
Los \qty{0.4}{L/min} de equivalente de oxígeno que faltan los aporta la
\omterm{def:g10:cell-metabolism:fermentation}{fermentación} láctica de los músculos. El ácido láctico se acumula; en
unos minutos los músculos arden y el corredor se ve obligado a bajar a
un ritmo que la respiración sola pueda cubrir.
\end{solution}

\begin{solution}{exo:g10:exercise-and-energy:14}
Las tandas cortas máximas van con la reserva fosforada y la
\omterm{def:g10:cell-metabolism:fermentation}{fermentación}, que no usan nada de grasa; la grasa solo la quema la
respiración, que tarda minutos en alcanzar su ritmo pleno y solo domina
en un ejercicio moderado y prolongado. Además, la energía total de unas
pocas tandas es pequeña frente a la de una hora de esfuerzo sostenido.
\end{solution}

\begin{solution}{exo:g10:exercise-and-energy:15}
Déficit $25000 - 8500 - 6000 = \qty{10500}{kJ}$, es decir unos
\qty{280}{g} de grasa. Comer durante la etapa mantiene alta la glucemia
para el cerebro y ahorra \omterm{def:g10:exercise-and-energy:glycogen}{glucógeno} para las subidas, donde la potencia
necesaria está muy por encima de lo que la grasa sola puede aportar; la
grasa cubre la parte constante del día.
\end{solution}

\begin{solution}{pb:g10:exercise-and-energy:1}
\textbf{1.} $4.0 \times 60 \times 42.2 \approx \qty{10100}{kJ}$.

\textbf{2.} $\qty{3}{h}\,\qty{30}{min} = \qty{12600}{s}$:
$\num{1.01e7}/12600 \approx \qty{800}{W}$.

\textbf{3.} $10100/20 \approx \qty{506}{L}$; en \qty{210}{min}, unos
\qty{2.4}{L/min}.

\textbf{4.} $2400/60 = \qty{40}{mL/min}$ por kilogramo, es decir
$40/52
\approx 77\%$ de su máximo.

\textbf{5.} $(2.4 - 0.2)/2.4 \approx 92\%$.

\textbf{6.} $380 \times 17 = \qty{6460}{kJ}$.

\textbf{7.} La carrera cuesta \qty{240}{kJ} por kilómetro: $6460/240
\approx \qty{27}{km}$.

\textbf{8.} Consumo de \omterm{def:g10:chemistry-of-life:families}{glúcidos} $0.8 \times 240 = \qty{192}{kJ/km}$:
$6460/192 \approx \qty{34}{km}$ --- el muro.

\textbf{9.} Energía de la grasa $0.2 \times 10100 \approx \qty{2020}{kJ}$, es decir $2020/38 \approx \qty{53}{g}$: alrededor del
0,6\% de sus \qty{9}{kg}.

\textbf{10.} La grasa solo la quema la respiración y a velocidad
limitada: puede aportar aproximadamente la mitad de la potencia que
exige su ritmo. Sin \omterm{def:g10:exercise-and-energy:glycogen}{glucógeno}, los músculos no fabrican ATP con
suficiente rapidez y el ritmo cae.

\textbf{11.} $60 \times 3.5 = \qty{210}{g}$, es decir $210 \times 17 \approx
\qty{3570}{kJ}$, que cubren $3570/192 \approx \qty{19}{km}$ de
consumo de \omterm{def:g10:chemistry-of-life:families}{glúcidos}.

\textbf{12.} \omterm{def:g10:chemistry-of-life:families}{Glúcidos} disponibles $6460 + 3570 = \qty{10030}{kJ}$:
$10030/192 \approx \qty{52}{km}$, más allá de la meta. Sí.

\textbf{13.} $+\qty{200}{g}$ de \omterm{def:g10:exercise-and-energy:glycogen}{glucógeno} $+ \qty{600}{g}$ de agua:
\qty{0.8}{kg}, que cuestan $0.8 \times 4.0 \approx \qty{3}{kJ}$ más por
kilómetro --- despreciable frente a 240.

\textbf{14.} \omterm{def:g10:exercise-and-energy:glycogen}{Glucógeno} $580 \times 17 = \qty{9860}{kJ}$ más la bebida:
unos \qty{13400}{kJ}; consumo de \omterm{def:g10:chemistry-of-life:families}{glúcidos} $0.55 \times 240 =
\qty{132}{kJ/km}$: más de \qty{100}{km}. El muro se ha ido mucho más
allá de la meta, con amplio margen.

\textbf{15.} El cerebro quema solo glucosa, que le suministra el
\omterm{def:g10:exercise-and-energy:glycogen}{glucógeno} del hígado. Cuando esa reserva se vacía, la glucemia cae, y
el cerebro señala la urgencia como mareo y fatiga abrumadora, tengan lo
que tengan aún los músculos.

\textbf{16.} $3.6 \times 55 \times 42.2 \approx \qty{8360}{kJ}$;
$8360/20 = \qty{418}{L}$ en \qty{125}{min}: unos \qty{3.3}{L/min}.

\textbf{17.} $3340/55 \approx \qty{61}{mL/min}$ por kilogramo, es decir
$61/80 \approx 76\%$ --- la misma fracción que Nadia. Lo que difiere es
el techo, no la fracción que se usa de él.

\textbf{18.} La economía es el coste por kilogramo y por kilómetro.
Para un corredor de \qty{55}{kg}, $(4.0 - 3.6) \times 55 \times 42.2 \approx
\qty{930}{kJ}$ ahorrados en la carrera, alrededor de un 9\%.

\textbf{19.} \omterm{def:g10:exercise-and-energy:glycogen}{Glucógeno}, \qty{8500}{kJ}; \omterm{def:g10:chemistry-of-life:families}{glúcidos} necesarios $0.7 \times
8360 \approx \qty{5850}{kJ}$: sí, las reservas bastan sin beber.

\textbf{20.} Sin preparación, el \omterm{def:g10:exercise-and-energy:glycogen}{glucógeno} de Nadia se agota hacia el
kilómetro 34; la carga de \omterm{def:g10:chemistry-of-life:families}{glúcidos} antes de la carrera, beber \omterm{def:g10:chemistry-of-life:families}{glúcidos}
durante ella y un entrenamiento que aumente la proporción de grasa
quemada llevan cada uno el muro más allá de la línea de meta.
\end{solution}
